\section{Protótipo - Extração de Landmarks}

Protótipo simplificado da extração de landmarks utilizando MediaPipe:

\begin{lstlisting}[language=Python]
import cv2
import mediapipe as mp

# Inicializa MediaPipe
mp_holistic = mp.solutions.holistic
holistic = mp_holistic.Holistic()

# Captura vídeo
cap = cv2.VideoCapture(0)

while True:
    ret, frame = cap.read()
    if not ret:
        break

    # Processa frame
    rgb_frame = cv2.cvtColor(frame, cv2.COLOR_BGR2RGB)
    results = holistic.process(rgb_frame)

    # Extrai landmarks das mãos
    if results.left_hand_landmarks:
        for lm in results.left_hand_landmarks.landmark:
            print(f"Mão esquerda: x={lm.x}, y={lm.y}, z={lm.z}")

    if results.right_hand_landmarks:
        for lm in results.right_hand_landmarks.landmark:
            print(f"Mão direita: x={lm.x}, y={lm.y}, z={lm.z}")

    cv2.imshow('Frame', frame)
    if cv2.waitKey(1) & 0xFF == ord('q'):
        break

cap.release()
cv2.destroyAllWindows()
\end{lstlisting}

\section{Protótipo - Treinamento do Modelo LSTM}

Protótipo simplificado do modelo LSTM para classificação:

\begin{lstlisting}[language=Python]
import tensorflow as tf
from tensorflow.keras.models import Sequential
from tensorflow.keras.layers import LSTM, Dense

# Constrói modelo simples
model = Sequential([
    LSTM(64, activation='relu', input_shape=(30, 543)),
    Dense(32, activation='relu'),
    Dense(3, activation='softmax')  # 3 classes: biblioteca, prova, secretaria
])

model.compile(optimizer='adam', loss='categorical_crossentropy',
              metrics=['accuracy'])

# Simula treinamento com dados aleatórios
import numpy as np
X_train = np.random.rand(100, 30, 543)  # 100 sequências
y_train = np.random.randint(0, 3, (100, 3))

model.fit(X_train, y_train, epochs=10, batch_size=32)
\end{lstlisting}

\section{Protótipo - Interface em Tempo Real}

Protótipo simplificado da interface com predição:

\begin{lstlisting}[language=Python]
import cv2
import mediapipe as mp
import numpy as np

# Carrega modelo treinado
model = tf.keras.models.load_model('modelo_libras.h5')

# Inicializa captura
cap = cv2.VideoCapture(0)
class_names = ['biblioteca', 'prova', 'secretaria']

while True:
    ret, frame = cap.read()
    if not ret:
        break

    # Extrai landmarks
    rgb_frame = cv2.cvtColor(frame, cv2.COLOR_BGR2RGB)
    results = holistic.process(rgb_frame)

    # Simula predição
    landmarks = np.random.rand(1, 30, 543)
    prediction = model.predict(landmarks)
    class_idx = np.argmax(prediction[0])
    confidence = prediction[0][class_idx]

    # Exibe resultado
    if confidence > 0.7:
        text = f"{class_names[class_idx]}: {confidence:.2f}"
        cv2.putText(frame, text, (50, 50),
                   cv2.FONT_HERSHEY_SIMPLEX, 1, (0, 255, 0), 2)

    cv2.imshow('LibrasVision', frame)
    if cv2.waitKey(1) & 0xFF == ord('q'):
        break

cap.release()
cv2.destroyAllWindows()
\end{lstlisting}

\section{Descrição do Dataset de Treinamento}

O dataset foi coletado seguindo os critérios:

\begin{itemize}
    \item \textbf{Total de sinais coletados:} 150 vídeos (50 por classe)
    \item \textbf{Usuários participantes:} 5 indivíduos surdos fluentes em Libras
    \item \textbf{Duração média por sinal:} 2-3 segundos (60-90 quadros a 30 FPS)
    \item \textbf{Resolução dos vídeos:} 640x480 pixels
    \item \textbf{Condições de captura:} Iluminação natural e artificial, fundo neutro
    \item \textbf{Variações incluídas:} Diferentes velocidades de execução, ângulos e distâncias da câmera
\end{itemize}

\section{Consentimento dos Participantes}

Este projeto seguiu os protocolos éticos necessários para pesquisa envolvendo seres humanos. Todos os participantes assinaram Termo de Consentimento Livre e Esclarecido (TCLE) autorizando o uso de seus vídeos para fins de pesquisa, com garantia de anonimato e privacidade.

Os dados foram armazenados de forma segura e apenas os landmarks (coordenadas abstratas) foram utilizados para treinamento, não armazenando os vídeos originais no modelo final.

\section{Requisitos do Sistema}

Para executar o LibrasVision em ambiente de produção:

\begin{table}[htb]
    \centering
    \caption{Requisitos mínimos do sistema}
    \label{tab:requisitos_sistema}
    \begin{tabular}{|l|l|}
        \hline
        \textbf{Componente} & \textbf{Especificação Mínima} \\
        \hline
        Processador & Intel i5 ou equivalente \\
        RAM & 4 GB \\
        Câmera & Webcam 720p \\
        Python & Versão 3.8+ \\
        Sistema Operacional & Windows 10+, Linux, macOS \\
        \hline
    \end{tabular}
\end{table}